% GENERATED FILE. Edit canonical lessons, then run npm run generate:books.
% canonical-source-hash: fnv1a64:4d78d58e212c14db
% canonical-lessons: ZH-C119-shubao, ZH-W119-bao-wrap, ZH-C119-xiangzi, ZH-C119-hezi, ZH-C119-daizi, ZH-C119-xingli

\chapter{Bags and Boxes}
\label{ch:zh-bags-and-boxes}
\hlchaptermodality{119}

\begin{chapteropening}
\textbf{By the end of this chapter:} \emph{I can name a schoolbag, a suitcase, a box, a carrier bag and luggage.}

Complete the last lesson of chapter 119: I can name a schoolbag, a suitcase, a box, a carrier bag and luggage.
\end{chapteropening}

\section[\texorpdfstring{\zh{书包} (shūbāo)}{书包 (shūbāo)}]{\zh{书包} (shūbāo) — schoolbag}

\label{lesson:ZH-C119-shubao}



\begin{quote}
Before the new one: say the Mandarin for text message, then the Mandarin for to receive.
\end{quote}

\subsection*{You'll want to know: \zh{书包}}
\textbf{\zh{书包}} — \emph{shūbāo} — \textquotedblleft{}schoolbag\textquotedblright{}.

\begin{quote}
\textbf{\zh{我的书包里有两本书}} \emph{wǒ de shūbāo lǐ yǒu liǎng běn shū} — There are two books in my schoolbag.
\end{quote}

The new character, in pieces: \zh{包} — a wrapping stroke that bends round from the top, and inside it three strokes: a turn, a short bar, and a vertical that bends right along the bottom and hooks up. The next lesson puts it in your hand.

\begin{cousinweb}
\textbf{\zh{书包}} (\emph{shūbāo}) is built from \zh{书} 'book' and \zh{包} 'to wrap', both inherited from Middle Chinese, \emph{syo} and \emph{paew} (Guangyun, in Baxter's transcription): a wrapper for books.
\end{cousinweb}

\begin{grammarlens}[title={What you've built}]
One more word for everyday things.
\end{grammarlens}

\subsection*{Your turn}
\begin{itemize}
\raggedright
  \item \emph{Say it:} \emph{shūbāo}
  \item \emph{Say it:} \emph{shūbāo}, once more
  \item \emph{Say it:} \emph{wǒ de shūbāo lǐ yǒu liǎng běn shū} — There are two books in my schoolbag.
  \item \emph{Recall:} say the Mandarin for text message, then the Mandarin for to receive, then say \emph{shūbāo} again
\end{itemize}

\begin{tcolorbox}[breakable,colback=teal!4,colframe=teal!35!black,title={Before you move on}]
What does \zh{书包} mean? (\textquotedblleft{}Schoolbag\textquotedblright{}.) Say it once more.
\end{tcolorbox}
\section[\texorpdfstring{\zh{包} (bāo)}{包 (bāo)}]{\zh{包} — a wrapping stroke, and three strokes inside}

\label{lesson:ZH-W119-bao-wrap}



\begin{quote}
The word first, then the character inside it.

\begin{itemize}
\raggedright
  \item \emph{Recall:} say \emph{shūbāo} — \textbf{R1}, one lesson back, before it decays
\end{itemize}

\begin{quote}
\zh{书} \zh{也}
\end{quote}

\emph{Read it:} each one aloud, with nothing said first — \emph{shū}, \emph{yě}
\end{quote}

\subsection*{Script — a wrapping stroke, and three strokes inside}
\hlblockfigure{\detokenize{figures/ZH-W119-bao-wrap-filmstrip.pdf}}{How \zh{包} is written, stroke by stroke}

\begin{quote}
\zh{包}
\end{quote}

\textbf{The wrapping piece, two strokes}: a short left-falling stroke, then a stroke across that turns down the right side and hooks.

\textbf{The inside, three strokes}: a stroke across that turns down, a short horizontal under it, and a vertical that bends right along the bottom and hooks up, out past the wrapping, as the last stroke of \zh{也} does.

2 + 3 = five.

\subsection*{Your turn}
\begin{itemize}
\raggedright
  \item \emph{Write it:} \zh{包}, all five, with the model in front of you
  \item \emph{Cover:} write it again from memory, and count the pieces, not the strokes
\end{itemize}

\begin{tcolorbox}[breakable,colback=teal!4,colframe=teal!35!black,title={Before you move on}]
How many strokes in \zh{包}? (\textbf{Five}.) Which piece comes first, the wrapping or the inside? (\textbf{The wrapping}.) Which character you know ends with the same stroke? (\textbf{\zh{也}}.)

Source: \href{https://raw.githubusercontent.com/chanind/hanzi-writer-data/68d10a4b21150cae5e1ebbd223eed289cf32d90c/data/\%E5\%8C\%85.json}{Hanzi Writer Data \zh{包}, pinned stroke paths and medians}.
\end{tcolorbox}
\section[xiāngzi]{xiāngzi — box; suitcase; trunk}

\label{lesson:ZH-C119-xiangzi}



\begin{quote}
Before the new one: say the Mandarin for to receive, then the Mandarin for schoolbag.
\end{quote}

\subsection*{You'll want to know: xiāngzi}
\emph{xiāngzi} — \textquotedblleft{}box; suitcase; trunk\textquotedblright{}.

\begin{quote}
\emph{zhège xiāngzi tài dà le} — This suitcase is too big.
\end{quote}

Its characters come later; for now it is a word for the ear.

\begin{cousinweb}
\emph{xiāngzi} is built from xiāng and the noun ending zi, both inherited from Middle Chinese, \emph{sjang} and \emph{tsiX} (Guangyun, in Baxter's transcription); the dictionary glosses the first as a box or a hamper.
\end{cousinweb}

\begin{grammarlens}[title={What you've built}]
One more word for everyday things.
\end{grammarlens}

\subsection*{Your turn}
\begin{itemize}
\raggedright
  \item \emph{Say it:} \emph{xiāngzi}
  \item \emph{Say it:} \emph{xiāngzi}, once more
  \item \emph{Say it:} \emph{zhège xiāngzi tài dà le} — This suitcase is too big.
  \item \emph{Recall:} say the Mandarin for to receive, then the Mandarin for schoolbag, then say \emph{xiāngzi} again
\end{itemize}

\begin{tcolorbox}[breakable,colback=teal!4,colframe=teal!35!black,title={Before you move on}]
What does xiāngzi mean? (\textquotedblleft{}Box; suitcase; trunk\textquotedblright{}.) Say it once more.
\end{tcolorbox}
\section[hézi]{hézi — box; small case}

\label{lesson:ZH-C119-hezi}



\begin{quote}
Before the new one: say the Mandarin for schoolbag, then the Mandarin for suitcase.
\end{quote}

\subsection*{You'll want to know: hézi}
\emph{hézi} — \textquotedblleft{}box; small case\textquotedblright{}.

\begin{quote}
\emph{hézi lǐ yǒu shénme} — What's in the box?
\end{quote}

Its characters come later; for now it is a word for the ear.

\begin{cousinweb}
\emph{hézi} is built from hé and the noun ending zi, both inherited from Middle Chinese, \emph{hop} and \emph{tsiX} (Guangyun, in Baxter's transcription); the dictionary glosses the first as a dish with a lid.
\end{cousinweb}

\begin{grammarlens}[title={What you've built}]
One more word for everyday things.
\end{grammarlens}

\subsection*{Your turn}
\begin{itemize}
\raggedright
  \item \emph{Say it:} \emph{hézi}
  \item \emph{Say it:} \emph{hézi}, once more
  \item \emph{Say it:} \emph{hézi lǐ yǒu shénme} — What's in the box?
  \item \emph{Recall:} say the Mandarin for schoolbag, then the Mandarin for suitcase, then say \emph{hézi} again
\end{itemize}

\begin{tcolorbox}[breakable,colback=teal!4,colframe=teal!35!black,title={Before you move on}]
What does hézi mean? (\textquotedblleft{}Box; small case\textquotedblright{}.) Say it once more.
\end{tcolorbox}
\section[dàizi]{dàizi — bag; sack}

\label{lesson:ZH-C119-daizi}



\begin{quote}
Before the new one: say the Mandarin for suitcase, then the Mandarin for small box.
\end{quote}

\subsection*{You'll want to know: dàizi}
\emph{dàizi} — \textquotedblleft{}bag; sack\textquotedblright{}.

\begin{quote}
\emph{nǐ xūyào dàizi ma} — Do you need a bag?
\end{quote}

Its characters come later; for now it is a word for the ear.

\begin{cousinweb}
\emph{dàizi} is built from dài and the noun ending zi, both inherited from Middle Chinese, \emph{dojH} and \emph{tsiX} (Guangyun, in Baxter's transcription); the dictionary glosses the first as a kind of sack.
\end{cousinweb}

\begin{grammarlens}[title={What you've built}]
One more word for everyday things.
\end{grammarlens}

\subsection*{Your turn}
\begin{itemize}
\raggedright
  \item \emph{Say it:} \emph{dàizi}
  \item \emph{Say it:} \emph{dàizi}, once more
  \item \emph{Say it:} \emph{nǐ xūyào dàizi ma} — Do you need a bag?
  \item \emph{Recall:} say the Mandarin for suitcase, then the Mandarin for small box, then say \emph{dàizi} again
\end{itemize}

\begin{tcolorbox}[breakable,colback=teal!4,colframe=teal!35!black,title={Before you move on}]
What does dàizi mean? (\textquotedblleft{}Bag; sack\textquotedblright{}.) Say it once more.
\end{tcolorbox}
\section[xíngli]{xíngli — luggage}

\label{lesson:ZH-C119-xingli}



\begin{quote}
Before the new one: say the Mandarin for small box, then the Mandarin for carrier bag.
\end{quote}

\subsection*{You'll want to know: xíngli}
\emph{xíngli} — \textquotedblleft{}luggage\textquotedblright{}.

\begin{quote}
\emph{nǐ de xíngli zài nǎr} — Where's your luggage?
\end{quote}

Its characters come later; for now it is a word for the ear.

\begin{cousinweb}
\emph{xíngli} is built from xíng 'to go' and lǐ 'plum', both inherited from Middle Chinese, \emph{haeng} and \emph{liX} (Guangyun, in Baxter's transcription). The dictionary lists the pair itself under the second, but how 'plum' came into a word for luggage is not settled.
\end{cousinweb}

\begin{grammarlens}[title={What you've built}]
One more word for everyday things.
\end{grammarlens}

\subsection*{Your turn}
\begin{itemize}
\raggedright
  \item \emph{Say it:} \emph{xíngli}
  \item \emph{Say it:} \emph{xíngli}, once more
  \item \emph{Say it:} \emph{nǐ de xíngli zài nǎr} — Where's your luggage?
  \item \emph{Recall:} say the Mandarin for small box, then the Mandarin for carrier bag, then say \emph{xíngli} again
\end{itemize}

\begin{tcolorbox}[breakable,colback=teal!4,colframe=teal!35!black,title={Before you move on}]
What does xíngli mean? (\textquotedblleft{}Luggage\textquotedblright{}.) Say it once more.
\end{tcolorbox}
